\section{从真实本地大模型推理引擎开始}

\subsection{第三册的主线不是玩具模型}

AI 推理最容易被讲成一组名词：tensor、operator、layer、quantization、GEMM、SIMD、attention、KV cache、runtime、scheduler。名词本身没有错，但如果没有机器账本，它们很快会变成框架文档的索引。第三册的主线必须更具体：写出一台本地推理引擎，让它能加载开源 7B 左右 decoder-only Transformer 权重，把 prompt 分词成 token id，完成 prefill 和 decode，在 CPU 后端先跑通，再逐步优化，并且让同一套 runtime 能挂接 GPU 后端。

一个真实请求不是“调用模型”四个字，而是下面这条链路：

\begin{lstlisting}[numbers=none,caption={真实本地大模型推理请求的主链路}]
prompt text
  -> tokenizer.encode()
  -> token embedding
  -> repeated transformer blocks
       RMSNorm
       Q/K/V projection
       RoPE
       attention with KV cache
       output projection
       MLP gate/up/down projection
  -> final norm and lm_head
  -> sampler
  -> next token
  -> append token and repeat decode
\end{lstlisting}

这条链路里有两种完全不同的性能形态。prefill 阶段一次处理 prompt 中的一串 token，很多投影和 MLP 子层更像 GEMM；decode 阶段每次只追加一个 token，许多地方退化成 GEMV 或小 batch GEMM，同时还要读取越来越长的 KV cache。若不区分这两种阶段，就无法解释为什么同一个模型在长 prompt、短回复、长上下文和 batch 服务下会表现不同。

\begin{keyidea}
第三册的贯穿项目是一台真实本地大模型推理引擎。最小 MLP 只是第一块显微镜切片，用来手工检查张量、矩阵乘、量化和 oracle；它不是本册目标，也不是最终架构。
\end{keyidea}

\subsection{先画清楚引擎边界}

一台本地大模型引擎至少有七个边界。第一，模型文件边界：权重、shape、dtype、量化参数、层数、hidden size、head 数、词表大小和 tokenizer 配置必须能被一致读取。第二，runtime 边界：session 保存当前 token 序列、position、KV cache、临时 activation buffer 和采样状态。第三，算子边界：RMSNorm、matmul、RoPE、attention、softmax、MLP、sampler 都要有独立语义合同。第四，后端边界：同一个算子可以由 CPU kernel 或 GPU kernel 执行，但输入输出合同不能变化。第五，内存边界：权重、packed weight、activation、KV cache 和后端工作区有不同生命周期。第六，调度边界：prefill、decode、batch、线程池、GPU stream 和拷贝同步不能混成一团。第七，证据边界：正确性、误差、性能、峰值内存和对标对象必须写进报告。

\begin{systemdiagram}{真实推理引擎的七个边界}
\begin{pdfdiagram}[x=1cm,y=1cm]
  \node[book accent node,text width=2.25cm] (model) at (0,2.15) {模型文件\\weights/config};
  \node[book teal node,text width=2.25cm] (tok) at (3.0,2.15) {tokenizer\\text -> ids};
  \node[book purple node,text width=2.25cm] (rt) at (6.0,2.15) {runtime\\session/state};
  \node[book amber node,text width=2.25cm] (graph) at (9.0,2.15) {decoder graph\\operators};

  \node[book navy node,text width=2.25cm] (mem) at (0,0.25) {内存计划\\arena/KV};
  \node[book teal node,text width=2.25cm] (backend) at (3.0,0.25) {后端合同\\CPU/GPU};
  \node[book red node,text width=2.25cm] (sched) at (6.0,0.25) {调度\\prefill/decode};
  \node[book purple node,text width=2.25cm] (evidence) at (9.0,0.25) {证据\\oracle/bench};

  \draw[book strong arrow] (model) -- (tok);
  \draw[book strong arrow] (tok) -- (rt);
  \draw[book strong arrow] (rt) -- (graph);
  \draw[book weak arrow] (model) -- (mem);
  \draw[book weak arrow] (graph) -- (backend);
  \draw[book weak arrow] (rt) -- (sched);
  \draw[book strong arrow] (backend) -- (evidence);
  \draw[book strong arrow] (sched) -- (evidence);
  \draw[book dashed arrow] (mem) -- (backend);

  \node[book label,text width=9.8cm] at (4.5,-1.25) {本图要先证明一件事：大模型推理引擎不是一个 kernel，而是模型表示、状态、内存、后端、调度和证据共同组成的系统。};
\end{pdfdiagram}
\diagramnote{七个边界分别约束文件、文本、运行状态、算子语义、内存生命周期、执行后端和验收证据。}
\end{systemdiagram}

这张图也解释了为什么第三册不能只讲 CPU。CPU 是第一条完整落地路径，因为它最适合把地址流、cache、SIMD、多线程、NUMA 和反汇编讲清楚；GPU 是第二条后端路径，因为真实本地 7B 推理经常需要显存带宽和矩阵吞吐。两条路径共享同一份模型语义和算子合同，只是在 kernel 实现、内存驻留、调度同步和 profiler 证据上不同。

\subsection{用最小切片建立显微镜}

虽然主线是真实大模型引擎，第二章仍然从一个极小 MLP 切片开始。原因不是降低目标，而是建立显微镜。一个只有线性层、bias、activation、量化权重和 argmax 的模型，可以把推理引擎中最重要的几件事压缩到一页纸上：

\begin{lstlisting}[numbers=none,caption={最小切片的推理语义}]
hidden = relu(input * W0 + b0)
logits = hidden * W1 + b1
class_id = argmax(logits)
\end{lstlisting}

这段伪代码看起来只是在做矩阵和向量运算，但机器实际看到的是输入向量的连续字节、权重矩阵的布局、bias 的加载、乘加累加器、激活函数的分支或 \instruction{max} 指令、输出缓冲的写入，以及最终 \code{argmax} 的比较序列。若模型使用 int8 权重量化，核心问题又会变成：8 位权重怎样解码，输入 scale 和权重 scale 怎样组合，累加器为什么通常使用 32 位，重新量化时怎样舍入和饱和。

这些问题在 7B 模型里不会消失，只会被放大。MLP 切片中的 \code{input * W0} 对应 Transformer 里的 Q/K/V projection、output projection 和 MLP gate/up/down projection；\code{hidden} 对应 activation buffer；量化权重对应数十亿参数的压缩格式；oracle 对应每层输出、logits 和生成 token 的误差验证。后续内容会不断把显微镜下的局部机制接回真实大模型引擎。

\subsection{张量不是数组别名}

\term{tensor}{张量} 可以先理解为“带形状的多维数据视图”，但这个说法还不够。对推理引擎来说，张量至少包含五类信息：

\begin{itemize}
  \item 元素类型：例如 \code{float}、\code{std::int8_t}、\code{std::int32_t}、\code{bf16}、\code{fp16}。
  \item 形状：每个维度有多少元素，例如 \code{[tokens, hidden]}、\code{[heads, seq, head_dim]}。
  \item 步长：相邻维度在内存中跨多少元素或字节。
  \item 量化元数据：scale、zero point、group size、per-channel 或 per-group 参数。
  \item 所有权和驻留位置：CPU 内存、GPU 显存、mmap 权重、临时 arena 或 KV cache。
\end{itemize}

连续的一维数组只是最简单的张量。二维矩阵若按行主序存储，访问 \code{A[i,j]} 常接近：

\begin{lstlisting}[numbers=none,caption={行主序矩阵的地址计算}]
address = base + (i * stride_row + j) * sizeof(element)
\end{lstlisting}

若矩阵被转置视图、切片视图或 padding 后的视图引用，\code{stride_row} 和 \code{stride_col} 就不再等于直觉中的宽度。算子如果假设所有输入都连续，就可能在正确性或性能上出错；算子如果完全支持任意 stride，又可能让热循环多出复杂地址计算。工程实现必须明确边界：哪些 kernel 要求 contiguous，哪些路径负责转换布局，转换成本是否被计入测量。

在大模型推理里，张量还会变成状态。KV cache 不是普通临时张量：它跨 decode 步骤持续存在，每生成一个 token 都要向每层的 K/V 区域追加一行，并在后续 attention 中被重新读取。它的 shape、stride、容量、有效长度和后端驻留位置都属于模型正确性与性能合同。

\subsection{算子合同：先定义结果，再谈 kernel}

\term{operator}{算子} 是有输入合同和输出合同的计算单元。以矩阵乘为例，语义可以写成：

\begin{lstlisting}[numbers=none,caption={矩阵乘 reference 语义}]
for m in 0..M:
  for n in 0..N:
    acc = 0
    for k in 0..K:
      acc += A[m, k] * B[k, n]
    C[m, n] = acc
\end{lstlisting}

这段 reference 不追求快，它定义正确结果。优化 kernel 可以换循环顺序、分块、展开、向量化、预打包权重，也可以使用 int8 乘加指令或 GPU tensor core；但只要算子合同没有变，输出就必须被 oracle 接受。浮点路径可能用误差界比较，整数量化路径可能用反量化后的误差界比较，生成路径还要检查 logits 排名、采样前概率和最终 token 序列是否在合同内。

\begin{warningbox}
不要先写“高性能 kernel”再倒推语义。AI 推理里很多 bug 并不崩溃，而是悄悄改变 logits、概率、排名或边界 token。没有 reference 和误差 oracle，速度数字没有意义。
\end{warningbox}

\subsection{量化从哪里来}

\term{quantization}{量化} 的基本动机不是神秘技巧，而是机器账本约束。一个 7B 模型如果权重按 float32 保存，仅权重就接近 28GB；按 fp16/bf16 保存约 14GB；按 int8 保存约 7GB；若使用 int4 或分组量化，权重字节还能继续下降，但会引入更复杂的 scale、zero point、group metadata 和反量化开销。权重变小可以降低内存容量压力、cache/TLB 压力、内存带宽压力和 GPU 显存压力，但量化不是免费压缩，它引入表示误差和重新缩放成本。

最常见的线性量化可以写成：

\begin{lstlisting}[numbers=none,caption={线性量化和反量化}]
q = clamp(round(x / scale) + zero_point, qmin, qmax)
x_approx = (q - zero_point) * scale
\end{lstlisting}

这里 \code{scale} 决定实数间隔，\code{zero_point} 让某个整数值表示实数零。若张量范围很大，scale 变大，小差异会被吞掉；若范围估计太窄，极端值会被 clamp。量化校准要解决的正是这个问题：用什么数据估计范围，用 per-tensor、per-channel 还是 per-group scale，怎样在误差、内存和速度之间取舍。

整数矩阵乘常把 int8 输入和 int8 权重乘出中间结果，再累加到 int32：

\begin{lstlisting}[numbers=none,caption={int8 dot product 的累加器}]
int32 acc = 0
for k in 0..K:
  acc += int32(a_i8[k] - za) * int32(w_i8[k] - zw)
out = requantize(acc, scale_a, scale_w, scale_out, zero_out)
\end{lstlisting}

累加器不能仍然用 int8，因为 \code{K} 稍大就会溢出。重新量化也不是简单右移：真实 scale 往往不是 2 的幂，需要乘法、舍入、饱和和边界处理。第三册后续会把这些步骤分别落到 C++ 类型、CPU/GPU 指令形态和测试 oracle。

\subsection{从算子到推理引擎}

单个算子正确还不等于推理引擎正确。decoder-only 推理会把多个算子按层重复连接起来：上一层输出成为下一层输入，中间 activation 可以被复用，权重可能预打包，某些算子可以融合，某些布局转换必须插入，KV cache 要跨 token 保留。推理引擎至少要管理五类状态：

\begin{itemize}
  \item 模型状态：权重、shape、量化参数、层配置、tokenizer 配置、算子拓扑。
  \item 运行状态：token ids、position、logits、临时 activation、KV cache、采样器状态。
  \item 后端状态：CPU feature、线程池、packed weight、GPU buffer、stream、kernel registry。
  \item 证据状态：reference 输出、误差报告、benchmark 配置、二进制身份、对标框架版本。
  \item 资源状态：内存峰值、cache 工作集、线程亲和性、NUMA 边界、显存占用和传输边界。
\end{itemize}

真实引擎的执行可以写成：

\begin{lstlisting}[numbers=none,caption={本地大模型推理引擎执行流程}]
load model metadata, tokenizer and weights
validate tensor shapes, dtypes and quantization parameters
prepare backend-specific packed weights or device buffers
session = initialize KV cache, activation arena and sampler

tokens = tokenizer.encode(prompt)
run prefill(tokens, session)

while not stopped:
    logits = run decode_one_token(session)
    next_id = sampler.sample(logits)
    session.append(next_id)
return tokenizer.decode(generated_ids)
\end{lstlisting}

这段流程故意没有隐藏工程边界。加载模型时要校验文件；选择 kernel 时要检查 dtype、layout 和 backend capability；复用 activation buffer 时要保证生命周期不重叠；更新 KV cache 时要保证层号、head、position 和 batch 维度一致；记录 timing 时要把 prefill tokens/s、decode tokens/s、首 token 延迟、内存峰值和对标对象分开。

\subsection{学完这一册能做出的业务原型}

把本节这些边界连起来，读者能做出的不是完整商业平台，而是一个严肃的本地推理原型。它应该具备：

\begin{itemize}
  \item 加载一个受限 decoder-only 模型或模型切片。
  \item 验证 tokenizer、manifest、shape、dtype、量化参数。
  \item 支持 prefill/decode，维护 session 和 KV cache。
  \item 输出流式 token、finish reason 和 usage。
  \item 支持至少一个 CPU reference path 和一个优化 path。
  \item 用 oracle 比较单算子、层输出或 logits。
  \item 用 profiler 报告 load、prepare、prefill、decode、sample。
  \item 在容量不足、context 超限、backend 不支持时给出明确诊断。
\end{itemize}

这已经足以支撑本地助手、RAG 问答、批量摘要或代码补全的原型。真正上线还需要安全策略、模型评测、权限、日志脱敏、监控、弹性调度和产品层交互，但底层 AI infra 能力已经闭合：输入可验证，状态可管理，后端可替换，误差可检测，性能可解释。

\begin{keyidea}
第三册要教的不是“调用一个大模型 API”，而是实现本地推理系统的关键骨架：模型合同、张量地址、量化数值、KV 状态、后端计划、oracle 和 profiler。
\end{keyidea}

\subsection{本节收束}

第三册从真实 7B 级本地推理引擎立题，再用小 MLP 切片建立显微镜。张量把内存、形状、量化和驻留位置连接起来；算子合同把语义和测试连接起来；量化把表示范围、内存容量和机器成本连接起来；KV cache 把单次算子变成跨 token 状态；后端合同把 CPU SIMD、多线程和 GPU kernel 放进同一个 runtime；性能报告把每一步优化放到当前基线和成熟框架旁边比较。

下一章不会急着跳到最快 GEMM 或某个 SIMD 指令，而是先把一层矩阵向量乘拆成 shape、stride、地址公式、float32 内层循环、基本块流图、抽象汇编账本和 int8 账本。只要这条最小路径能讲清，后面的 packing、向量化、量化校准、算子融合、attention、KV cache、CPU/GPU 后端和运行时调度才有稳定落点。
